% Драфт секции 1.6: Интерпретация внутренних представлений в ИНС
% Источники: NeuroAI_ideas.pdf, Neuro_for_AI.pdf, neurotrain/main.tex
% Статус: черновик для ревью

\section{Интерпретация внутренних представлений в искусственных нейронных сетях}\label{sec:lit_ann_repr}

% [старая связка, закомментирована]
% \link{Связка: предыдущие секции были посвящены методам анализа активности биологических нейронных сетей. Однако задача интерпретации внутренних представлений актуальна и для искусственных нейронных сетей. Более того, глубокие нейронные сети демонстрируют замечательные параллели с биологическими системами в организации нейронного кода, что позволяет переносить аналитические подходы между областями.}

Рассмотренные выше методы~--- снижение размерности, оценка размерности, сетевой анализ, анализ селективности~--- разрабатывались для активности биологических нейронных сетей. Однако задача интерпретации внутренних представлений в равной мере актуальна и для искусственных сетей, а глубокие сети обнаруживают выраженные параллели с мозгом в организации кода~--- от селективности отдельных элементов до низкоразмерной структуры популяционной активности. Эти параллели делают перенос аналитического инструментария между нейронаукой и машинным обучением не только возможным, но и продуктивным; ему посвящён настоящий раздел.

\subsection{Вычислительная роль отдельных элементов}\label{sec:ann_individual}

Поскольку искусственные нейронные сети представляют собой вычислительные графы с известной структурой, их можно <<препарировать>>~--- систематически анализировать вклад каждого элемента в работу сети~\cite{bau2017network, Olah2017}. Ранние работы показали, что отдельные нейроны глубоких свёрточных сетей формируют иерархию специализаций, аналогичную иерархии зрительной коры: нейроны первых слоёв реагируют на простые признаки (границы, текстуры), а нейроны глубоких слоёв~--- на высокоуровневые объекты и понятия~\cite{erhan2009visualizing, Olah2017}. Примечательно, что подобная селективность может быть мультимодальной~--- так, в модели CLIP были обнаружены нейроны, одновременно реагирующие на изображения и текстовые описания одного понятия~\cite{goh2021multimodal}.

Важным вопросом является степень <<моносемантичности>> отдельных нейронов~--- отвечает ли каждый нейрон за один конкретный признак, или его функция определяется совместно с другими элементами сети. Работы по механистической интерпретируемости показали, что в малых моделях можно выделить моносемантические единицы~\cite{Anthropic2023monosemanticity}, однако в больших сетях преобладают \textit{полисемантические} нейроны, реагирующие на несвязанные совокупности входов~\cite{Templeton2024scaling}. Для выделения интерпретируемых признаков из полисемантических нейронов применяются разрежённые автокодировщики (\textit{sparse autoencoders}, SAE), проецирующие пространство активаций в более высокоразмерное, но разрежённое пространство отдельных <<фич>>~\cite{Templeton2024scaling}.

Эта картина перекликается с наблюдениями в нейробиологии: наряду с высокоселективными <<бабушкиными нейронами>>~\cite{concepts} и клетками места~\cite{o1971hippocampus}, большинство нейронов мозга демонстрируют смешанную селективность~--- нелинейную зависимость от нескольких переменных задачи одновременно (см. раздел~\ref{sec:mixed_selectivity}). Распределение силы специализаций в обоих случаях подчиняется логнормальному закону: небольшая доля элементов обладает ярко выраженной селективностью, тогда как основная масса вносит слабый, но коллективно значимый вклад~\cite{Buzsaki2014}.

\subsection{Популяционная динамика активности искусственных нейронов}\label{sec:ann_population}

Ограничения индивидуального анализа~--- полисемантичность элементов, суперпозиция признаков, комбинаторный рост числа возможных взаимодействий~--- мотивируют переход к популяционному уровню описания, аналогично тому, как это происходит при анализе биологических данных (см. раздел~\ref{sec:population_coding}).

Согласно \textit{гипотезе линейных представлений}~\cite{Park2023LRH}, внутренние понятия в глубоких сетях кодируются как направления в пространстве активаций. На основе этого принципа развивается область \textit{representation engineering}~\cite{Zou2023RepE}~--- подход <<сверху вниз>>, который анализирует глобальную структуру пространства представлений, не требуя доступа к отдельным нейронам или цепочкам. Было показано, что таким образом можно не только классифицировать внутренние состояния модели, но и целенаправленно управлять её поведением~--- например, усиливать или ослаблять <<направление честности>> в языковых моделях~\cite{Zou2023RepE}.

Однако не все внутренние представления линейны: в пространстве активаций были обнаружены круговые структуры (дни недели, месяцы) и более сложные нелинейные многообразия~\cite{Engels2024nonlinear}. Внутренняя размерность представлений в глубоких сетях меняется от слоя к слою, достигая минимума в средних слоях, что соответствует максимальному сжатию информации~\cite{Ansuini2019}. Этот результат перекликается с принципом <<информационного бутылочного горлышка>> и наблюдениями об изменении размерности нейронных представлений в мозге (см. раздел~\ref{sec:dimred_methods}).

Глубокие сети демонстрируют устойчивость, характерную для популяционного кода: из ResNet можно удалить целые слои без существенной потери качества~\cite{Veit2016residual}, а языковые модели обладают механизмами <<самовосстановления>> при повреждении отдельных компонентов~\cite{McGrath2023hydra}. Эти свойства согласуются с распределённым характером нейронного кодирования и затрудняют интерпретацию на уровне отдельных элементов.

Таким образом, параллели между организацией представлений в искусственных нейронных сетях и в мозге~--- селективность элементов, низкоразмерная структура популяционной активности, устойчивость к повреждениям~--- делают возможным и продуктивным перенос аналитических инструментов между областями. В Главах~\ref{sec:collective} и~\ref{sec:individual} настоящей работы эта возможность реализуется на практике: методы снижения размерности и взаимной информации, разработанные для нейрофизиологических данных, применяются к анализу активности искусственных РНС (раздел~\ref{sec:rnn}), а задача синтеза оптимальных стимулов, описанная в следующем разделе, решается одними и теми же методами для биологических и искусственных нейронов.
